\section{Dynamics: motor lag and the state ODE}
\label{sec:dynamics}

The kinematic map of \S\ref{sec:kinematics} treats wheel commands as instantaneously producing the desired body twist. Real motors have inertia, controller dynamics, and current limits; their realised speed lags the command. We fold all of this into a single first-order lag on the body velocities.

\subsection{Commanded body twist}

The wheel command pair $u = (v_L, v_R)$ produces a commanded body twist via the forward kinematic map \eqref{eq:fwd-kin}:
\begin{equation}
    v_{\text{cmd}} = \half(v_L + v_R),
    \qquad
    \omega_{\text{cmd}} = \frac{v_R - v_L}{\trackw}.
    \label{eq:body-cmd}
\end{equation}
``Commanded'' to distinguish from the realised body velocities $(v, \omega)$ that appear in the state.

\subsection{Motor lag}

The realised body velocities track the commanded values through a first-order lag with time constant $\taum$:
\begin{align}
    \dot v       &\;=\; \frac{1}{\taum}\,\bigl(v_{\text{cmd}} - v\bigr),
    \label{eq:vdot} \\
    \dot \omega  &\;=\; \frac{1}{\taum}\,\bigl(\omega_{\text{cmd}} - \omega\bigr).
    \label{eq:omegadot}
\end{align}
The same time constant is used for both channels. We could carry separate $\taum^{v}$ and $\taum^{\omega}$ if a motor identification campaign warranted it, but in practice the difference is absorbed into the per-episode domain randomisation in \texttt{training/}.

\begin{remark}[Why first-order, not second-order?]
A second-order lag would track an additional state per channel (acceleration), capturing the motor's torque dynamics more faithfully. For the wheel masses and drive sizing we expect, the second-order modes are well above the policy's control bandwidth and a first-order fit is adequate. The model can be promoted to second-order later without disturbing anything outside this section.
\end{remark}

\subsection{Full state ODE}

Stack the position, heading, and body twist into a single state vector $\bx \in \R^{5}$:
\begin{equation}
    \bx \;=\; \begin{pmatrix} p_x & p_y & \theta & v & \omega \end{pmatrix}^{\!\top}.
\end{equation}
Combining the unicycle kinematics \eqref{eq:unicycle} with the motor lag \eqref{eq:vdot}--\eqref{eq:omegadot} and substituting \eqref{eq:body-cmd}:
\begin{equation}
    \dot \bx \;=\; f(\bx,\, u) \;=\;
    \begin{pmatrix}
        v \cos\theta \\[2pt]
        v \sin\theta \\[2pt]
        \omega \\[4pt]
        \dfrac{1}{\taum}\!\left(\half(v_L + v_R) \,-\, v\right) \\[6pt]
        \dfrac{1}{\taum}\!\left(\dfrac{v_R - v_L}{\trackw} \,-\, \omega\right)
    \end{pmatrix}.
    \label{eq:state-ode}
\end{equation}
This is the equation \texttt{DiffDriveDynamics::continuous} in \texttt{core/dynamics.hpp} implements verbatim. The integrator (\texttt{rk4\_step} in \texttt{core/integrators.hpp}) treats $u$ as zero-order-hold over a step of length $\Delta t$ and advances $\bx$ accordingly.

\subsection{Why 5-D, not 7-D?}

There are three reasonable parameterisations:

\begin{description}
    \item[3-D kinematic.] State $(p_x, p_y, \theta)$, input $(v, \omega)$. No actuator dynamics --- fine for kinematic planners but useless for tracking control.
    \item[5-D body-lag.] State $(p_x, p_y, \theta, v, \omega)$, input $(v_L, v_R)$. Single first-order lag at the body level. \emph{This document.}
    \item[7-D wheel-lag.] State $(p_x, p_y, \theta, \dot\phi_L, \dot\phi_R, \ldots)$, input $(v_L, v_R)$. Per-wheel lag, body twist computed downstream.
\end{description}

We chose the middle option: it captures the dominant transient (motor lag) without paying for two redundant per-wheel states, and it keeps the embedded MPC problem small enough that linearised solvers like TinyMPC stay in their sweet spot. If a future task demonstrates per-wheel lag matters --- e.g., uneven loading or one motor saturating before the other --- the model expands to 7-D, but evidence-driven, not preemptively.

\subsection{Slip}

The model assumes pure rolling. The kinematic map \eqref{eq:fwd-kin} carries no slip term, and the motor lag operates on the resulting body velocities directly. Domain randomisation on $\trackw$, $\taum$, and the per-wheel gain absorbs much of what slip would contribute in practice. If a deployed policy starts exploiting non-physical kinematic behaviour --- e.g., turning faster than $\omega_{\text{cmd}}$ permits --- a velocity-dependent slip term such as
\begin{equation}
    \omega_{\text{achieved}} \;=\; \omega_{\text{cmd}}\,(1 - \alpha\,|v|)
\end{equation}
can be folded into \eqref{eq:body-cmd}, but only as a corrective response to observed exploit, not as an a-priori complication.
